\section{群の作用}


\subsection{軌道}
群は作用を考えてからが本番。
\begin{prop}
有限群$G$が集合$X$に作用しているとする. $x\in X$の固定部分群を$G_x$, 軌道を$O_x$で表すとき,
\ben
\item \bolm{|G_x||O_x| = |G|}
\item
\een
\end{prop}
\prf (1): $a\in O_x$とする. $O_x = \set{x,g_2x, \dots, g_n x}$とし, これらに重複がないとする($|O_x| = n$). 写像$\phi:G\to O_x$を$\phi(h) = hx$により定めると, $hx = g_i x\iff g_i^{-1}h \in  G_x$である. よって, 任意の$i$に対して $\phi^{-1}(g_ix) = g_i G_x$だから
\[G = \bigcup_{i=1}^{n} \phi^{-1}(g_i x) = \bigcup_{i=1}^{n}g_iG_x \]
から集合の元の個数を計算すれば得られる. \qed

\begin{eg}
位数$2m$ ($m$は奇数) の群$G$において, 位数2の元はSylow-2部分群を作る. Sylowの定理から位数2の元はすべて共役であることが分かる. よって, 位数2の元全体は一つの共役類$X$をなす(共役類上で位数は一定であることに注意). $X$には$G$が共役で作用する.  もしある位数2の元$x$が$\cyc{2k}$ に同型な部分群に含まれているとする. このとき, $x$の固定部分群$G_x$, すなわち, すなわち中心化部分群$C_{G}(x)$はその部分群を含んでいるので, $|X|$は$2k$の倍数であることが分かる.
\end{eg}







\subsection{cosetへの作用}

次の問題は覚えておくとよい。
\tbox{
\begin{egprob}
$G$を有限群, $p$を$G$の位数を割り切る最小の素数とする。任意の指数$p$の$G$の部分群は正規部分群であることを証明せよ。
\end{egprob}
}{title=H18東工大}
\begin{egans}
指数$p$の部分群$H$を取る. cosetの集合$G/H$を考え, 置換表現$\rho:G\to \mr{Aut}(G/H)$; $gH \xmapsto{g'} g'gH $を考える. $\mr{Aut}(G/H)$の位数は$p!$である.
\begin{claim}
$H = \ker{\rho}$である. とくに, $H$は正規部分群である.
\end{claim}
[証明] $[G:H]=p$, $[G:\ker{\rho}] = |\im{\rho}|$である. \aka{ $k\in \ker{\rho}$であるとき, とくに$kH = H$であることから$k\in H$なので, $\ker{\rho} \subset H$が従う.}(これを見落としがちである)  よって$[G:\ker{\rho}] = [G:H][H:\ker{\rho}] = p[H:\ker{\rho}]$であるから$[H:\ker{\rho}]=1$を示せばよい. そこでこれが$[H:\ker{\rho}] > 1$であったとしよう. いま,
\[\frac{1}{p}|\im{\rho}| = [H:\ker{\rho}]\]
が分かっている. $|\im{\rho}|$は$p!$の約数だから, $[H:\ker{\rho}]$は$(p-1)!$の約数である. よって仮定よりこの数は$p$未満の素因数$q$を持つ. $|G| = [G:H][H:\ker{\rho}][\ker{\rho}:\set{e}]$だから$|G|$の約数でもある. よって$|G|$は$q$で割り切れるがこれは$p$の最小性に反する. \qed
\end{egans}



ガロア理論と絡めて出題されるようなこともある。正直, 嫌な問題だと思う。
\begin{egprob} (H1-I-3)\label{prob:H1_I_3}
可換体$K$の5次拡大体$K(\theta)$で, $K(\theta)$を含む$K$の最小のGalois拡大体が$K$上の40次になるものは存在しないことを示せ。
\end{egprob}
\begin{screen}
やけに問題が一般的だと思わないだろうか？5次拡大体なんていろいろあるだろうに, このようなことが成り立つというのなら,  この現象が群論レベルの出来事である可能性が高い。
\end{screen}
\begin{ans}
$5$次対称群$\maf{S}_{5}$に指数$3$の部分群は存在するかを考えればよい。ここで使う非自明な事実は, $A_5$が単純群であるということである。\par
\tf{指数3の部分群は存在しない。}これを背理法によって示す。

$H$を$\maf{S}_{5}$の指数$3$(位数$40$)の部分群とする。$40$は$60$を割りきらないから, $H$は$A_{5}$の部分群ではない。よって$H$には奇置換が存在し, 奇置換と偶置換がそれぞれ$20$個ずつ存在する($\because$奇置換$\sigma$を一つ選び, $H\ni g\mapsto g\sigma\in H$という$H$の自己同型が$H$の奇置換と偶置換の全単射を与えるから)。$K:=H\cap A_5$とすると, $K$は位数$20$である。cosetの集合$X=A_{5}/K$を考え, $A_5$の$X$への作用を$\sigma K \mapsto (\tau\cdot \sigma )K$で定める。つまり, 群準同型$\phi:A_5\to \maf{S}_{3}$を
\[\phi(\tau) = (f_{\tau}: \sigma K \mapsto (\tau \sigma )K)\]
で定義する。このとき$\ker{\phi}$は$A_5$の正規部分群であるが, $A_5$の単純性より$\ker{\phi} = \{ 0 \}, A_5$となるしかないが, 写像$A_5\to S_3$が単射になることはありえないので$\ker{\phi} \neq \{ 0 \}$であり, 一方$\ker{\phi} = A_5$だと$\phi$は自明な準同型だが $\phi$の定義から明らかに不適である。よって矛盾が得られた。\qed
\end{ans}



\subsection{演習問題}
\subsubsection{Level B}

$S_5$に指数3の部分群が存在しないことを以前用いたが,  $S_4$には存在するというのも有名である. 次はそれが題材になっている.

\begin{prob} (H31東北大専門1)
$G=\maf{S}_{4}$の二つの元$\sigma = (1234)$, $\tau = (24)$で生成される部分群を$H$とする.
\ben
\item $H$が非可換であることを示せ.
\item $G/H$の元のそれぞれに対して代表元を一つずつ求めよ.
\item $K=\set{g\in G\mid \forall x\in g, (gx)H = xH}$は$G$の正規部分群であることを示せ.
\item $G/K\cong \maf{S}_{3}$を示せ.
\een
\end{prob}

\subsubsection{Hint・Answer}
{\small
\begin{probans}
(1): $\sigma\tau \sigma^{-1} = (\sigma(2)\sigma(4)) = (31)\neq \tau$より. \\
(2): $H$の位数が8であることを示す. まず, $G$の元に対する次の性質を考える.
\lefba{ ($\ast$)
$g\in G$を全単射$g:\set{1,2,3,4}\to \set{1,2,3,4}$とみなすとき, $g(\set{2,4}) = \set{2,4}$であるかまたは$g(\set{2,4}) = \set{1,3}$である.
}
$g_1,g_2,g\in G$が$(\ast)$を満たせば$g_1g_2$と$g^{-1}$もこれを満たすので, $(\ast)$を満たす元全体は部分群をなす. $\sigma,\tau$はこれを満たすので, $H$のすべての元はこれを満たす. また, $(\ast)$を満たす$g$は, 組み合わせ的に8通りしか存在しないので, $H$の位数は8の約数. $H$は$1,\sigma,\sigma^2,\sigma^3,\sigma\tau$の5つの元を含むので, $\sharp{H}=8$が分かる. $H$の元は$(\ast)$で特徴づけられるということが分かったので, $(12),(34)\notin H$である.
\lefba{ ここで $(12)H \neq (34)H$を示す. もし等号が成り立つなら$(12)\sigma(34) \in H$だが, これは$(23)$であり, $(\ast)$を満たさないので$H$に属さない. よって矛盾.
}
よって$G/H$は3つの元を持ち, $1,(12),(34)$で代表される. \\
(3): $K=\bigcap_{x\in G}x^{-1}Gx$と書けるので$h^{-1}Kh = K$はすぐ分かる. \\
(4): $K$はcosetへの作用$\phi:G\to \mr{Aut}(G/H) \cong \maf{S}_{3}$の核である. よって$G/K\cong \im{\phi}$.
\[\sigma^2 = (13)(24), \sigma\tau = (12)(34), \sigma^3\tau = (23)(14)\]
はサイクル分解で, 同じ型を持っているので$G$で共役である. (2)の代表元を用いると,
\[(12)\sigma^2(12) = (14)(23) = (34)\sigma2(34) \in H\]
なので, $\sigma^2 \in K$である. $K$は正規だから$K$は$\sigma^2$の共役元も含み, 位数4以上であることが分かる. 一方, $(12)\tau(12) = (14)\notin H$なので$\tau\notin K$であり, $\sharp{K} = 4$が確定する. よって$G/K$も$\maf{S}_{3}$も位数が6なので$G/K\to \maf{S}_{3}$は全単射準同型であり, 同型である. \qed
\end{probans}

}



\newpage
